\section{Related Work}
\label{sec:related}

\subsection{Flow-based intrusion detection for MQTT and IoT}

MQTT-IoT-IDS2020 was released with a baseline study comparing six classical
learners across packet, unidirectional-flow and bidirectional-flow abstractions
of the same captures~\cite{hindy2020mqttcase}. Its central observation is one
this paper returns to: general network attacks such as scans separate easily
from benign traffic, whereas MQTT-specific attacks resemble normal broker
exchanges and are harder to isolate. Subsequent work on the dataset has
largely reported higher aggregate accuracy rather than revisiting that
distinction. Khan et al.~\cite{khan2021mqttdl} apply deep models to the
MQTT-enabled setting; ensemble approaches have been proposed for MQTT
specifically~\cite{mqttensemble2025,mqttsecuring2023}, with feature selection
by metaheuristic search~\cite{mqttlightgbm2022}; attention-based sequence
models have been applied to MQTT in medical
deployments~\cite{mqttiomt2026}; and XMID-MQTT~\cite{xmidmqtt2025} adds
post-hoc explanation. MQTTset~\cite{vaccari2020mqttset} provides an
alternative MQTT capture, and broader IoT benchmarks such as
TON\_IoT~\cite{moustafa2021toniot} and CICIoT2023~\cite{neto2023ciciot} offer
larger and more heterogeneous traffic.

The base paper for this study~\cite{norouzi2026reliable} proposes GA-EBT, a
genetic algorithm that jointly selects features and tunes a bagging ensemble of
decision trees, and reports 99.9931\% and 99.997\% accuracy on IoT\_SDN-IDS and
MQTT-IoT-IDS2020 respectively. We reproduce its Algorithm~1 as faithfully as
the published description permits and discuss the comparison in
Section~\ref{sec:results}; one property of its experimental setup, identified in
Section~\ref{sec:dataset}, materially affects how its figure should be read.

\subsection{Graph neural networks for intrusion detection}

Two surveys map this area~\cite{bilot2024survey,gnnids2026survey}. The
dominant formulation represents hosts as nodes and flows as edges, which makes
intrusion detection an edge-classification problem.
E-GraphSAGE~\cite{lo2022egraphsage} established this formulation for IoT by
extending GraphSAGE~\cite{hamilton2017graphsage} so that edge attributes
participate in neighbourhood aggregation, and it remains the closest published
comparator to the model proposed here; we therefore include it as a
first-class baseline rather than as a variant.
Anomal-E~\cite{caville2022anomale} removes the dependence on labels by
learning edge embeddings self-supervised and applying anomaly scoring on top.
Pujol-Perich et al.~\cite{pujolperich2022unveiling} make robustness the
explicit target, arguing that structural representations should degrade more
gracefully than flow-vector models under traffic that differs from training.
More recent work enhances edge semantics through self-supervised
objectives~\cite{semantic2026edge}.

On the architectural side, attention over graphs originates with
GAT~\cite{velickovic2018gat}, whose limitations under static attention were
analysed by Brody et al.~\cite{brody2022gatv2}. Admitting edge features into
the attention computation is not new: EGAT~\cite{chen2021egat} computes
attention from a concatenation that includes a projection of the edge feature
vector, which is the same placement we use. We therefore make no novelty claim
for the mechanism. What is not established in the literature is how that
placement behaves for flow-level intrusion detection when the evaluation is
constructed so that endpoint identity cannot be memorised, and that is the gap
this paper addresses.

\subsection{Evaluation methodology}

A parallel literature argues that security machine learning is systematically
over-reported. Sommer and Paxson~\cite{sommer2010outside} set out why intrusion
detection resists the closed-world assumptions of standard supervised
evaluation. Arp et al.~\cite{arp2024pitfalls} enumerate recurring pitfalls,
among them sampling bias, spurious correlations and inappropriate baselines,
all three of which bear directly on graph-based detectors evaluated on
laboratory testbeds. TESSERACT~\cite{pendlebury2019tesseract} shows for malware
classification that eliminating experimental bias can remove most of a reported
improvement. Dataset construction itself has been found defective:
Engelen et al.~\cite{engelen2021troubleshooting} and
Lanvin et al.~\cite{lanvin2022errors} document labelling and generation errors
in CICIDS2017 large enough to change which methods appear to win.
Geirhos et al.~\cite{geirhos2020shortcut} give the general account of shortcut
learning that we appeal to when interpreting endpoint identity as a shortcut.
On metrics, Chicco and Jurman~\cite{chicco2020mcc} argue for the Matthews
correlation coefficient over accuracy and F1 under imbalance, and
Saito and Rehmsmeier~\cite{saito2015prcurve} for precision-recall over ROC;
we follow both.

\subsection{Positioning}

Table~\ref{tab:related} summarises the comparison. The pattern it shows is
that graph-based intrusion detection is evaluated almost exclusively under
random splits of a single capture, that endpoint-identity baselines are absent,
and that where robustness is addressed it is addressed by perturbing traffic
rather than by holding out devices. The gap we act on is therefore not the
absence of an edge-aware attention mechanism, which
exists~\cite{chen2021egat}, but the absence of an evaluation under which its
contribution---or the contribution of any graph structure---can be measured
rather than assumed.

\begin{table*}[!t]
\caption{Representative prior work on graph-based and MQTT intrusion
detection. ``Endpoint holdout'' asks whether any reported protocol withholds
entire hosts from training; ``identity baseline'' whether a model using only
endpoint identity is reported as a comparator.}
\label{tab:related}
\centering
\small
\begin{tabular}{@{}llllcccl@{}}
\toprule
Study & Dataset(s) & Model & Graph & Edge feat. & Endpoint holdout &
Identity baseline & Principal limitation \\
\midrule
Hindy et al.~\cite{hindy2020mqttcase} & MQTT-IoT-IDS2020 & 6 classical & no &
n/a & no & no & random split; accuracy-led \\
Norouzi \& G\"urka\c{s}-Ayd{\i}n~\cite{norouzi2026reliable} & MQTT-IoT-IDS2020,
IoT\_SDN-IDS & GA + bagged trees & no & n/a & no & no &
saturated metric; subset of release \\
Khan et al.~\cite{khan2021mqttdl} & MQTT-IoT-IDS2020 & CNN/LSTM & no & n/a &
no & no & random split \\
Lo et al.~\cite{lo2022egraphsage} & ToN-IoT, BoT-IoT & GraphSAGE & yes & yes &
no & no & random split; no identity control \\
Caville et al.~\cite{caville2022anomale} & NF-UQ-NIDS & GraphSAGE + DGI &
yes & yes & no & no & unsupervised; random split \\
Pujol-Perich et al.~\cite{pujolperich2022unveiling} & CIC-IDS2017 & custom MPNN
& yes & yes & no & no & robustness via traffic perturbation \\
Chen \& Chen~\cite{chen2021egat} & citation/molecule graphs & EGAT & yes & yes
& n/a & n/a & not an intrusion-detection study \\
Bilot et al.~\cite{bilot2024survey} & survey & --- & yes & varies & rarely &
no & identifies evaluation gap, does not close it \\
\midrule
This work & MQTT-IoT-IDS2020 (uni- and bi-flow) & E-GAT-IDS + 11 baselines &
yes & yes & \textbf{yes} & \textbf{yes} & single dataset; simulated testbed \\
\bottomrule
\end{tabular}
\end{table*}

\subsection{Research gap}

Existing flow-based intrusion detection generally treats network observations
as independent feature vectors, and existing graph-based intrusion detection
demonstrates that relational representations can be applied to the same task at
comparable or better aggregate accuracy. Neither line establishes how much of
the graph's apparent contribution survives an evaluation in which endpoint
identity is not predictive, because no reported protocol withholds endpoints
and no reported comparison includes an endpoint-identity baseline. On a
simulated testbed with a fixed attacker address this omission is consequential,
and MQTT-IoT-IDS2020 is such a testbed.
